\raggedbottom
\section{Scoring objective, assumptions, and evidence status}
\label{app:objective}

This section distinguishes the documented score, our raw-rank diagnostic, and
claims that require assumptions about opponents or unplayed actions.

For the focal agent, these distinctions explain why the design could prefer a
lower expected monetary payoff, why agreement or purchase counts alone do not
establish policy quality, and why most strategy revisions are described
observationally rather than causally. They also bound the claim made for the two
raw-rank proxy decisions in the core: the bargaining Bob CDF raise in
Table~\ref{tab:policy} and the hidden one-round seller selector in
Eq.~\eqref{eq:hidden-seller}. Both were motivated by the scoring objective, but
neither is a proven optimizer of the private adjusted score.

\paragraph{Score model.}
A game has public configuration $c$, role $r$, history $H_t$, own observed type
$\theta_i$, possibly hidden opponent type $\theta_{-i}$, terminal payoff $Y$,
and opponent $o$. Let $\mathcal I_{i,t}$ contain every input visible to the
policy at time $t$, including the complete request, any disclosed opponent
identity, and the fixed policy configuration. The platform converts $Y$ to a
same-configuration, same-role percentile $q$, then to an opponent-adjusted percentile
$\widetilde q$ used in Eq.~\eqref{eq:rating} \citep{gleecompetition2026}.
The public documentation does not give the adjustment function or a tie rule.
We therefore use the exact-equality empirical midrank only as a local
diagnostic. For a fixed nonempty reference multiset
$D=\{z_1,\ldots,z_m\}$, $m\geq1$, define
\begin{equation}
Q_D(y)=\frac{1}{m}\sum_{j=1}^{m}
\left[\mathbf 1\{z_j<y\}+\frac12\mathbf 1\{z_j=y\}\right].
\label{eq:midrank}
\end{equation}
This is a probability-of-superiority statistic, since
$Q_D(y)=\Prb(Z<y)+\tfrac12\Prb(Z=y)$ for $Z$ sampled uniformly from $D$
\citep{mann1947test}. It is not asserted to be the organizer's tie convention.

\begin{proposition}[One-step ordering and terminal-offer identity]
Fix the current raw rating $R$, reference pool, opponent, adjustment, and the
same action-independent learning rate $\eta>0$ for one normally completed game.
For action $a$, let $\widetilde q(a)$ and $R_n^{\mathrm{raw}}(a)$ denote the
adjusted percentile and next raw rating under the outcome distribution induced
by choosing $a$, and define
$\Delta R^{\mathrm{raw}}(a)=R_n^{\mathrm{raw}}(a)-R$. Suppose the rating clamp
is inactive for every outcome with positive probability under either action.
Then, for any actions $a$ and $b$,
\[
\E[\Delta R^{\mathrm{raw}}(a)]\geq\E[\Delta R^{\mathrm{raw}}(b)]
\quad\Longleftrightarrow\quad
\E[\widetilde q(a)]\geq\E[\widetilde q(b)].
\]
Separately, in a normally completed terminal interaction with only acceptance
or rejection, if proposal $x$ is accepted with probability $A(x)$, pays $u(x)$
when accepted, and pays zero when rejected, then its expected local raw rank is
\begin{equation}
\E[Q_D(Y)\mid x]=Q_D(0)+A(x)\{Q_D(u(x))-Q_D(0)\}.
\label{eq:terminal-rank}
\end{equation}
\end{proposition}
\begin{proof}
Substitution in Eq.~\eqref{eq:rating} gives
$\E[\Delta R^{\mathrm{raw}}(a)]
=\eta\{2000+8000(\E[\widetilde q(a)]-1/2)-R\}$,
whose only action-dependent term has a positive coefficient. For
Eq.~\eqref{eq:terminal-rank}, condition on acceptance and rejection and collect
terms.
\end{proof}

The two statements use different ranks. Against $D=\{5\}$, a sure payoff of 10
has $(\E[Y],\E[Q_D(Y)])=(10,1)$, whereas an equal lottery over 0 and 21 has
$(10.5,1/2)$. Expected payoff and expected raw rank can therefore reverse, but
the example orders the private official score only if its adjustment preserves
that ordering. Also, $Q_D(0)$ need not be zero: it equals $1/2$ when every
reference payoff is zero. These qualifications make the implemented rules rank
proxies, not exact optimizers of the private platform score.

The one-step equivalence is not a leaderboard theorem. It excludes special
handling of abandonment or dropped games, activity penalties, decay, a binding
clamp, and changes in the reference pool or opponent adjustment. Display
shrinkage at fixed game count preserves the ordering, but an evolving field and
future opponent sequence need not.

\paragraph{Hidden information and unplayed actions.}
Suppose the same policy-visible information state $\mathcal I_{i,t}$ arises
under two latent opponent types that imply different profitable actions. A
deterministic policy $\pi_i(\mathcal I_{i,t})$ must choose the same action at
time $t$; later actions can differ only after the observed histories diverge. A
prior or acceptance model can support a Bayesian
decision, but it is a modeling input, not a fact recovered from the visible
state. Logs reveal the response to the action actually played, not the terminal
outcome under an unplayed action and continuation policy. Identifying that
counterfactual requires action variation and either a known randomized
assignment mechanism with recoverable realized assignments or, for
observational multi-turn data, sequential conditional exchangeability. Both
routes require consistency, positive assignment probability for the compared
actions, and prespecified handling of missing outcomes, stopping, and
interference. Alternatively, an identified and correctly specified structural
response model can supply the counterfactual under its assumptions. This is why
the cost ceiling inferred from a hidden seller offer remains a behavioral
assumption and why the observed-leader and selected-trigger analyses in the core
are descriptive rather than counterfactual.

\paragraph{Evidence status.}
Table~\ref{tab:evidence-status} fixes the evidence terms used throughout the
paper. The second-place endpoint is therefore treated as realized-system
evidence, not as an estimate of the final policy or any one component.

\begin{table}[H]
\centering
\small
\caption{Evidence objects and the strongest claim each supports.}
\label{tab:evidence-status}
\begin{tabularx}{\linewidth}{@{}>{\raggedright\arraybackslash}p{2.85cm}YY@{}}
\toprule
Evidence object & Supports & Does not support \\
\midrule
Formal derivation & Consequences of stated score and game assumptions. & Field optimality or opponent adherence. \\
Returned-action audit & Agreement of recorded owner, gate, or plan fields with the normalized returned action where instrumented. & Unlogged eligibility or owner selection, complete process provenance, or outcome benefit. \\
Defined observational-corpus audit & Actions and outcomes observed under an explicit selection rule and denominator. & The response to an unplayed action or a component effect. \\
Authenticated replay or result join & Correspondence of local returned actions and outcomes with captured service replay moves and results. & Code-to-agent or agent-to-leaderboard identity binding, organizer correctness, or component attribution. \\
Competition endpoint & Realized close-time rating and placement of the evolving account. & The effect of the final policy, a revision, or any component. \\
Qualifying randomized analysis & An intent-to-treat contrast when the randomization mechanism and realized assignments are recoverable, arm probabilities are positive, treatment is isolated, every assignment is retained or follows a prespecified missing-outcome rule, and stopping and interference are addressed. & Effects outside enrollment or eligibility, or effects after post-treatment selection. \\
\bottomrule
\end{tabularx}
\end{table}

\FloatBarrier
\section{Final policy specification and precedence}
\label{app:policy-spec}

This section specifies the recorded intended bank corresponding to the two
audit windows in Section~\ref{sec:behavior}. Its object is the locally returned
action, not a component effect or proof that every live process loaded the
recorded state.

\subsection{Frozen scope and common action path}

The recorded composite slot launched the \texttt{hardliner} profile from source
commit \texttt{b54cb33}; a separate manifest
hashes 23 files from it. The profile supplied bargaining SPE weight 1 and clock
cap 60 before the effective overlay. Table~\ref{tab:policy-receipt} records
other non-source inputs. All recorded per-game hash arms were zero and the
custom randomized-probe router was off; unrecorded arms evaluate at frozen
source defaults. Thus, with LLM mode off, the policy was deterministic
conditional on the request, overlay, defaults, frozen assets, parser, cache
state, inherited environment, and source version.

\begin{table}[H]
\centering
\scriptsize
\caption{Recorded inputs to the final action function. Hashes identify the
intended bank, not the state loaded by each live process. Full SHA-256 values
appear in the accompanying manifest.}
\label{tab:policy-receipt}
\setlength{\tabcolsep}{3.0pt}
\renewcommand{\arraystretch}{0.94}
\begin{tabularx}{\linewidth}{@{}>{\raggedright\arraybackslash}p{2.2cm}Y@{}}
\toprule
Object & Recorded identifier or value \\
\midrule
Source & Commit \texttt{b54cb33ad9f4fe7da4c667262eee00ffaeac3ab6}; 23-file asset manifest. \\
Overlay receipts & Nonpersuasion \texttt{33c9a48838ac485f}; persuasion \texttt{9fa0fb417acb4c4e}. \\
Runtime defaults & Hidden bargaining $\delta=.90$ and cap 99; negotiation clock 10 and continuation probability $.65$. \\
Execution mode & LLM off; bargaining messages off; persuasion token style preserved. \\
\bottomrule
\end{tabularx}
\end{table}

An \emph{owner} is the first eligible rule that returns a candidate action;
family transforms may follow. The common coercer repairs the phase schema,
applies negotiation price postguards except in first-principles-owned states or
for exact recomputed deterministic wires, removes forbidden text, and truncates
text above 2,000 characters. That output is the \emph{returned action}; server
receipt is separate. An exception, invalid candidate, or unrepairable candidate
triggers a protocol-legal fallback and then a phase-shaped last resort. These
are not utility-safe: offers use the nearest legal equal bargaining split or
the negotiator's own value, decisions accept, and persuasion recommends or buys.

\begin{table}[H]
\centering
\scriptsize
\caption{Summary of priority in the recorded final policy. Within bargaining
and negotiation, later rows are reached only when earlier eligible rules return
no action. Persuasion rows are mutually exclusive action phases.}
\label{tab:policy-precedence}
\setlength{\tabcolsep}{3.0pt}
\renewcommand{\arraystretch}{0.96}
\begin{tabularx}{\linewidth}{@{}>{\raggedright\arraybackslash}p{1.25cm}>{\raggedright\arraybackslash}p{3.15cm}Y>{\raggedright\arraybackslash}p{2.65cm}@{}}
\toprule
Family & Applicable scope & Candidate action & Subsequent treatment \\
\midrule
Bargaining & P2 focal and current; silent; disclosed $T=12$; own $\delta_2=.95$ & Six-step share ladder; discounted terminal floor. & Exact-owner coercion on every owned offer. \\
Bargaining & All other supported states & Horizon planner, continuation floors, eligible Bob CDF raise. & Acceptance floor, pacing fallback, numeric repair. \\
Negotiation & Outer owner: seller offer, disclosed $T=1$ & Visible $b-.01$ or hidden support-grid ask. & First-principles-owned state; guard exempt. \\
Negotiation & V1 branch: complete buyer, disclosed $T=10$ & $s+.01$ wires; $.994$ preterminal hold threshold. & Exact recomputed wire; guard exempt. \\
Negotiation & Hidden $T=10$ buyer, terminal, usable bound & Largest consistent support cost, plus $.01$. & Protected wire; else fall through. \\
Negotiation & Remaining well-formed supported buyers & Fixed observed-leader price and decision paths. & Exact recomputed wires. \\
Negotiation & Sellers and unsupported or unowned states & Generic Boulware, stall, continuation, deadgame, endgame. & Text table, repair, then claim postguard. \\
Persuasion & Seller phase & HIGH yes; LOW by prior value or buy rate. & Disabled overlays, token encoding, coercion. \\
Persuasion & Buyer phase & V2 posterior and money gate. & Disabled overlays, credibility reset, V1 decline veto, coercion. \\
\bottomrule
\end{tabularx}
\end{table}

\subsection{Bargaining}
\label{app:policy-bargaining}

\paragraph{Narrow silent Bob owner.}
Eligibility requires the focal and current player to be P2, a valid bargaining
phase with correct proposer and parity, \texttt{horizon\_known=True}, numeric
$T=12$, numeric $\delta_2=.95$, \texttt{messages\_allowed=False}, and $M>.01$.
It does not inspect Alice's discount or the \texttt{complete\_information}
flag. On Bob's six even-round proposal opportunities, it requests Bob payoffs
\[
(.61M,.68M,.74M,.82M,.91M,M-.01).
\]
On odd decision round $r$, it accepts Bob payoff $b_r$ exactly when
Eq.~\eqref{eq:bargaining-floor} holds. For rounds 2--10, let $g=s_rM$ for the
listed share $s_r$: if $g$ lies within $10^{-9}\max(M,1)$ of an integer, Python
ties-to-even rounding is used; otherwise $g$ is preserved. Alice receives the
exact residual. The round-12 $M-.01$ wire remains exact. Every independently
verified owned offer bypasses generic rounding. This is internal discipline,
not an observable commitment or equilibrium claim.

\paragraph{Generic owner.}
Let $r$ be the current round. Missing own discount gives $\delta_i=1$; an
absent opponent-discount key gives the planner $\delta_j=.90$. A horizon is
treated as disclosed only when \texttt{horizon\_known} is not false and
\texttt{max\_rounds} is present. Set $T$ to that cap or otherwise 99, and
$k=\max(1,T-r+1)$. With the recursion $f_i$ in
Section~\ref{sec:strategy}, the pre-transform aspiration is
\begin{equation}
a_0=\min\{M,\max[f_i(k)M,.57M,C_i]\},\qquad
C_i=\begin{cases}0,&k=1,\\
\delta_i[1-f_j(k-1)]M,&k>1.
\end{cases}
\label{eq:banked-barg-aspiration}
\end{equation}
The $.57M$ term is an intermediate floor, not a returned-offer invariant. The
active realistic-continuation estimator has ceiling
$S_i=\delta_i f_i(k-1)M$, replaced by $.57M\delta_i$ when
$\delta_j\geq1$ and $k>2$. With fewer than two opponent offers it returns
$S_i$. Otherwise, let $o_0,o_1$ be the first and last gains offered to us,
$g=(o_1-o_0)/(r_1-r_0)$ with denominator at least one, and
$h_{\max}=\min(5,k-1)$. It returns
\begin{equation}
R_i=\begin{cases}
\min(S_i,\delta_i o_1),&g\leq0,\\
\min\{S_i,\max_{1\leq h\leq h_{\max}}(o_1+hg)\delta_i^h\},&g>0.
\end{cases}
\label{eq:banked-barg-realistic}
\end{equation}
With a hidden horizon, the concession clock is
\begin{equation}
H_i=\begin{cases}
\max\{4,\min[\lceil\log(.05)/\log\delta_i\rceil,60]\},&0<\delta_i<1,\\
60,&\delta_i\geq1,
\end{cases}
\label{eq:banked-barg-clock}
\end{equation}
which equals $14,29,59,60$ on the supported discount grid. For
$e=\min\{1,\max[0,(r-1)/H_i]\}$ and $L$ the greater of $R_i$ and the best
observed opponent offer, the hidden-horizon target is
$L+(a_0-L)(1-e)^2$ when $a_0>L$.

Offer transforms run in source order. At $k=1$ the planner first sets $.75M$;
at $k=2$ it uses $\max\{C_i,\min(a,.62M)\}$. The opponent floor then caps our
share at $.61M$. On eligible Bob offers with $k>2$, a fitted local-rank search
may raise, but never lower, that result. For candidate Bob share $x$ it evaluates
\begin{equation}
J(x)=\widehat P_{\rm acc}(1-x,r,\delta_A^{\rm raw},d)\widehat Q(x\delta_i^{r-1})
+[1-\widehat P_{\rm acc}(1-x,r,\delta_A^{\rm raw},d)]\widehat V_{\rm cont},
\label{eq:banked-bob-search}
\end{equation}
where $d$ is Alice's latest retained share and a missing Alice discount is
passed as \texttt{None}, activating the model's hidden-discount indicator. If
$\ell$ is the larger of $.20$, the current target share, and Bob's best received
share, candidates are $\operatorname{round}(\ell+.01i,4)$ for
$i=0,\ldots,\operatorname{round}[(.95-\ell)/.01]$; an off-grid $\ell$ can put
the last point slightly above $.95$. Ties go upward. The result is used only
when $J(x)>\widehat V_{\rm cont}+.005$; otherwise the pre-model target remains.
The CDF key is $(M,T^{\rm field},\mathbf1\{T\text{ known}\},
\mathbf1\{\text{complete}\})$, with $T^{\rm field}=\texttt{None}$ when
withheld, and falls back to a pooled CDF. Missing model or CDF also leaves the
target unchanged. Role, discounts, message mode, and opponent adjustment are
absent from the CDF, so this is a raw-rank proxy, not the official optimizer.

At computed $k\leq1$, including the synthetic round-99 boundary, the generic
owner accepts exactly when current gain is positive. For $k>1$, it first rejects
nonpositive gain. Otherwise its threshold starts at $.98R_i$. For $k>2$ it is
at least $F_iM$, where
\[
F_i=\min\{.50,\delta_i(.05)/(1-\delta_i)\}\quad(\delta_i<1),
\qquad F_i=.50\quad(\delta_i=1).
\]
No equality repair was banked: binary floats make the nominal $.20$ and $.45$
floors at $\delta_i=.8,.9$ slightly larger, so exact $.20M$ or $.45M$ offers can
fail \texttt{>=}. In the intended reconstruction, absent
\texttt{BARG\_CLIFF\_GUARD} is false; if the threshold still rejects a positive
hidden-horizon offer at $r\geq H_i$, the pacing fallback accepts. The configured
high-share delay child was unreachable because its deterministic discount-aware
no-regression parent was off. Generic integer-pot offers use Python rounding on
our gain and the exact residual; noninteger pots use cents with residual repair.

\subsection{Negotiation}
\label{app:policy-negotiation}

Write seller cost $s$, buyer value $b$, price $p$, and scale $B$. The outer
one-round seller owner runs before the V1 whole-policy router; it normally
shadows V1's duplicate one-round branch. It asks $b-.01$ when $b$ is visible,
using the midpoint only if $b-s\leq.02$. With hidden $b$, costs
$.8B,1.0B,1.2B,1.5B$ map respectively to
$1.2B-.01,1.2B-.01,1.5B-.01,2B$. The first three values maximize the declared
pooled seed-midrank proxy in Eq.~\eqref{eq:hidden-seller}; the last is a
dead-cell convention. V1's remaining action-changing branch is a
complete-information $T=10$ buyer: it offers or counters at $s+.01$ (or the
midpoint when $b-s\leq.02$), accepts before terminal only when its share of
$b-s$ is at least $.994$, and accepts any strictly profitable terminal offer.
A later priority branch handles a hidden round-10 buyer offer. Let $A_s$ be the
deduplicated prior seller offer, counter, and last-offer prices strictly below
$b$. If nonempty, it sets
\begin{equation}
\hat s=\max\{z\in\{.8B,1.0B,1.2B,1.5B\}:z<b,\ z\leq\min A_s\}
\label{eq:banked-hidden-bound}
\end{equation}
when the set is nonempty, then bids $\hat s+.01$. This inference assumes sellers
do not propose below cost. Without a usable bound, the state falls through.

\paragraph{Remaining buyers.}
The observed-leader table requires a game ID, on-grid own value, valid offer or
decision phase, explicit Boolean information and horizon markers, and either a
known cap of 1 or 10 or an unknown horizon with missing \texttt{max\_rounds}.
Unsupported or malformed buyers fall to the generic stack. An owned buyer
counts prior distinct own price wires and holds the last element after a
sequence ends. Complete unknown-horizon prices claim buyer shares
$(.960,.898,.836,.774,.712)$ of visible surplus. The compiled cap-10 sequence
$(.960,.898,.836,.774,.774)$ is shadowed by the structural owner above. Hidden
price paths in units of $B$ are
\begin{align*}
.8 &: (.360000,.471508,.569815,.652831,.716938,.752000),\\
1.0 &: (.830000,.830000,.830000,.830000),\\
1.2 &: (.860000,.860000,.860000,1.030000,1.030000),\\
1.5 &: (.905000,.905000,1.075000,1.075000,1.245000,1.245000,1.245000,1.245000).
\end{align*}
For complete information, share $\rho$ returns price $b-\rho(b-s)$. Prices are
clipped to $[0,b-.01]$. A true terminal receiver accepts exactly
when $b-p>0$; otherwise it accepts a profitable seller ask no larger than the
path's final price, rejects with the current path price, and in hidden unknown
horizons walks after eight distinct own prices. Message-enabled counters say
``Above what I can justify. Here is where I can go.'' This is an observed-leader
imitation rule, not recovered private code or an equilibrium. The recorded
buyer anchor and learned buyer acceptance rows are shadowed on these owned
paths.

\paragraph{Remaining sellers.}
In supported generic seller cells, a real cap has $T=10$ and
$R=\max(1,T-r+1)$; otherwise $R=\max(1,11-r)$. The policy uses
$e=1-(R-1)/10$, clipped to $.85$ only when uncapped. Let $h$ be the maximum
finite opponent buyer price in offer, counter, and last-offer records. Treating
$h$ as a lower bound on buyer value assumes buyer individual rationality. The
anchor is $a=\max(4s,h)$ when $h$ exists and $4s$ otherwise. For a positive
visible ZOPA $b>s$, set $z=s+.39(b-s)$; otherwise set $z=s+.02(a-s)$, with
$z\geq1.15s$ for a hidden $1.5B$ seller. The power-2 target is
$t=a+(z-a)e^2$. In a positive visible ZOPA it is then bounded by
\begin{equation}
t\leftarrow\min\{\max(t,s),\ b-.39(b-s)\}.
\label{eq:banked-nego-visible-bound}
\end{equation}

A stall exists when the last three opponent \texttt{history[*].offer} prices
have range at most $.01\max_j|p_j|$. A profitable stalled price is matched. On
an unprofitable stall, the seller parks at
$(.9B,1.4B,1.4B,\varnothing)$ for costs $(.8B,1.0B,1.2B,1.5B)$; if no park
exists it repeats the last own offer when present, otherwise retaining $t$.
After at least four opponent prices, the
$1.5B$ deadgame branch replaces $t$ by
$s+B\max[.10(1-e),.02]$ when the inferred bound is below $s+.02B$.
For continuation, write seller surplus $u(p)=p-s$. With fewer than two opponent
price records, $\widehat u=\max[0,u(t)]$. Otherwise a nonpositive first-to-last
concession rate gives $\widehat u=\max[0,u(p_{\rm last})]$; a positive rate is
projected for $1,\ldots,\min(4,R-1)$ steps and the largest nonnegative surplus
is retained. In both cases $\widehat u$ is capped by
$\max[0,u(a)]$, and $V_c=.65(1-1/R)\widehat u$.

The decision order is: true-terminal positive-gain acceptance; unprofitable
late uncapped deadgame walk after a detected stall; nonterminal visible-span
veto below $.48$; configured $.60$ final-option
comparison; hidden acceptance table; stall acceptance; $V_c$ acceptance; then
a role-improving counter. The frozen learned ask maps are empty, so the enabled
table never changes $t$ and acts only through hidden seller acceptance. Its
thresholds by elapsed quartile are
$(.498408369,.436318357,.034463609,.226859056)B$.
The $.60$ option is shadowed by the earlier $.48$ veto on its complete-information
scope, and the configured $.80$ ultimatum writer is also shadowed on supported
returned paths. Configured posterior and curve-pricing writers are likewise
action-shadowed on supported returned paths. The $.40$ span invariant is reached,
but is output-dominated by the $.48$ acceptance veto and later $.80$ claim floor.

The seller text table requires messages, a game ID, seller role, explicit
Boolean information and horizon markers, a supported cap, a price-bearing offer
or rejection, and well-formed prior seller offer records. Hidden cap 1 is
supported but shadowed by the outer owner; complete cap 1 is unsupported. Prior
seller price count 0 selects ``That is my price for this one. I would rather
close now than trade rounds \texttt{-{}-} tell me if it works.'' Counts 1--3
select ``Below my line, but here is a number I can actually sign.'' Counts 4--9
select ``I have moved as far as this is worth to me. Let us close here.'' At 10
or more, existing text is preserved. During coercion, after schema
normalization, generic complete-information seller prices claim exactly $.95$
of visible surplus on the first outgoing wire and at least $.80$ later, clamped
to $[s,b]$; message cleanup follows. Hidden sellers have no visible-span
postguard, and first-principles states and exact buyer wires bypass it.

\subsection{Persuasion}
\label{app:policy-persuasion}

\paragraph{Seller.}
Exact quality string \texttt{high} is always recommended. A buyer high value is
treated as known solely when its field is nonmissing. Otherwise, for LOW
quality, the seller recommends iff filtered history is empty or the cumulative
truthy buy rate is at least $.5$. With visible values, it recommends LOW iff prior value
$pv+(1-p)u$ weakly exceeds price. Otherwise the banked honesty horizon 99 and
zero lie shading make it truthful throughout the 20-round game. The runtime
overlay overrides more aggressive wrapper defaults. In text mode, token style
and message preservation return literal ``yes'' or ``no''; optional seller
overlays were off.

\paragraph{Buyer posterior and money gate.}
The buyer models the seller's LOW-recommendation rate on
$\mathcal Q=\{0,.125,.25,.5,.75,1\}$ with HIGH-message error $\epsilon=.02$.
Its fitted prior depends on binary or text mode and on prior bin
$p<.45$, $.45\leq p<.65$, or $p\geq.65$. In that order, the six frozen vectors
are
\begin{align*}
w^0_{\rm bin,lo}&=(.23423,.22560,.17998,.11431,.07904,.16684),\\
w^0_{\rm bin,mid}&=(.18906,.20381,.16814,.11501,.12142,.20256),\\
w^0_{\rm bin,hi}&=(.12644,.12828,.12814,.14198,.19069,.28446),\\
w^0_{\rm text,lo}&=(.16190,.13936,.10638,.07602,.08400,.43234),\\
w^0_{\rm text,mid}&=(.12912,.10903,.08988,.07497,.12686,.47014),\\
w^0_{\rm text,hi}&=(.08032,.07913,.08491,.12060,.21062,.42441).
\end{align*}
The loader rechecks at most every ten seconds, retains its last valid cached
document after a later load failure, and uses uniform weights if no valid model
has loaded. It clips prior weights below $10^{-4}$, floors history likelihoods
at $10^{-12}$, and returns to uniform weights if normalization degenerates.
For a historical parsed recommendation $m\in\{\mathrm{yes},\mathrm{no}\}$ with purchased
quality $z$, likelihood under $q\in\mathcal Q$ is
\begin{align*}
\ell_q(H,m)&=p[.98\mathbf1\{m=\mathrm{yes}\}+.02\mathbf1\{m=\mathrm{no}\}],\\
\ell_q(L,m)&=(1-p)[q\mathbf1\{m=\mathrm{yes}\}+(1-q)\mathbf1\{m=\mathrm{no}\}].
\end{align*}
For an unpurchased item, quality is marginalized. The current message then
reweights the resulting $q$ posterior, and
\[
\Prb(H\mid\mathrm{yes},q)=\frac{.98p}{.98p+q(1-p)},\qquad
\Prb(H\mid\mathrm{no},q)=\frac{.02p}{.02p+(1-q)(1-p)}.
\]
The mixture gives $p_H$. Let $\beta=(P-u)/(v-u)$, or one when $v\leq u$.
The baseline buys iff $p_Hv+(1-p_H)u\geq P$, except that a tight cell
$\beta\geq.75$ is blocked when $p<.65$. All rank, confidence-bound,
selectivity, strike, audit, and alternative-parser arms were off.

The base text parser returns positive for exact yes, y, or buy and negative for
exact no, n, or pass. Otherwise a weighted language score above $.5$ is positive,
below $-.5$ is negative, and the middle region abstains. Its frozen phrase
tables use last-contrast splitting, weight pre-contrast text by $.25$, apply a
24-character negation reach, and suppress overlapping matches. The banked
baseline collapsed any string abstention, including an empty string, to
positive, so it was optimistic on ambiguity. Historical nonstring signals can
fall back to a Boolean recommendation field, while the current V2 baseline does
not inherit that field when \texttt{seller\_message} is absent. Nondictionary
history entries are dropped; bought entries with invalid quality are
marginalized. The machine-readable phrase tables in the frozen source are
normative for arbitrary text replay.

\paragraph{Credibility reset and final veto.}
The reset runs after the baseline only in a finite tight buyer state with prior
money value no greater than price, a strict current positive signal, and
posterior expected value strictly above price. Its separate anchored parser
grants credibility exactly when the latest strict decline occurs after the
latest purchased LOW item. Invalid credibility changes BUY to PASS. Valid
credibility can change PASS to BUY only when the sole blocker was the $.65$
tight-prior fence. With no prior strict decline, credibility is invalid.
Malformed relevant history makes the reset ineligible and leaves the incoming
V2 action unchanged. The final V1 veto then changes BUY to PASS in text mode
when its expanded parser finds a current decline, or in binary mode for Boolean
false and, only when the message is absent, a false recommendation field. The
aliases ``n'' and ``pass'' are not strict binary negatives. The veto never
creates a purchase. Later V2 and V3 vetoes were off. If persuasion coercion
cannot repair a candidate, its legal fallback is permissive: seller text becomes
``I recommend this product.'', binary seller becomes yes, and buyer becomes yes.

\subsection{Coverage and verification boundary}

The artifact verifier hashes 23 frozen Git blobs and their ordered checksum. It
also counts 81 flag rows and checks six recorded default or constant values, but
does not reconstruct and rehash the two effective overlay dictionaries. The
formula-to-code verifier covers representative branches while importing the
later mutable worktree, not an isolated banked checkout. No verifier executes
banked dispatch, owner selection, and coercion over every audited state.

The manifest omits the fleet launcher and supervisor, \path{run_agent.py},
\path{glee_agent/probes.py}, the effective overlay dictionaries, and inherited
environment. It therefore cannot prove which checkout, imports, cache, or
overlay each process used. In particular, absent \texttt{GLEE\_LLM\_FAMILIES}
is reconstructed at its source default of persuasion; an inherited negotiation
entry could rewrite negotiation text even with LLM mode off.

The authenticated 245-game join matches 2,288 returned actions to valid replay
moves, but does not recover every owner or pre-coercion candidate. Negotiation
contributes only 19 pre-close games and 190 actions (87 seller, 103 buyer), with
none in the added close window. In persuasion, stricter retrospective parsing
finds four V1 misses among 80 text-buyer decisions, all in one game; all four
purchases revealed LOW. Under Table~\ref{tab:evidence-status}, this is an
implementation specification, not a policy-effect estimate or complete
deployment receipt.
